\documentclass[12pt]{ctexart}
\usepackage{a4}
\usepackage{graphicx}
\usepackage{geometry}
\usepackage{amsmath}
\usepackage{mathtools}
\usepackage{amssymb}
\geometry{a4paper,left=2cm,right=2cm,top=2.5cm,bottom=2.5cm}

\title{几何基础作业11}
\author{无名}
\date{\today}

\begin{document}
\maketitle

\section{设$P(x_0,y_0)$是椭圆上一点，利用仿射变换求过点$P$的椭圆切线方程}

椭圆方程:
\[C:\frac{x^2}{a^2}+\frac{y^2}{b^2}=1\]

圆方程：
\[C':x'^2+y'^2=1\]

构造映射
\[f,\mathbb{R}^2 \to \mathbb{R}^2: \frac{x}{a}=x',\frac{y}{b}=y'\]
显然它是一个一一映射

那么椭圆$C$经过仿射变换后就是$C'$，此时$P(x_0,y_0)$在$C'$上对应点$P'(\frac{x_0}{a},\frac{y_0}{b})$

对于单位圆$C'$，切线方程为$x_0x+y_0y=1$，
那么$P'$的切线方程就是$\frac{x_0}{a}x+\frac{y_0}{b}y=1$

进行逆仿射变换，即得到原方程
\[\frac{x_0x}{a^2}+\frac{y_0y}{b^2}=1\]椭圆在$P(x_0,y_0)$切线方程

\newpage
\section{找到一个关于圆锥曲线（或圆）的配极对偶命题组}
\begin{figure}[htb]
    \begin{center}
        \includegraphics[width=0.5\linewidth]{01.png}
        \caption{}
    \end{center}
\end{figure}
\begin{figure}[htb]
    \begin{center}
        \includegraphics[width=0.5\linewidth]{02.png}
        \caption{}
    \end{center}
\end{figure}

由于圆锥曲线经过射影变换可变成圆，故考虑圆
\subsection{帕斯卡定理}
设六边形$A_1A_2A_3 A_4 A_5 A_6$内接于圆$O$，那么
\[A=A_4\bigcap A_5 \hspace{2em} B=A_5A_6\bigcap A_2A_3 \hspace{2em} C=A_1A_6\bigcap A_3A_4
\]
共线
\subsection{布利安桑定理}
设六边形$B_1B_2B_3B_4B_5B_6$的边都切与圆$O$，那么$B_1B_4$,$B_2B_5$,$B_3B_6$交于同一点
\subsection{帕斯卡证明布利安桑}
如图，由配极，$B_i$的极线是$A_iA_{i+1}$

\[A=A_1A_2\bigcap A_4A_5 \hspace{2em} B=A_5A_6\bigcap A_2A_3 \hspace{2em} C=A_1A_6\bigcap A_3A_4
\]
由帕斯卡，ABC共线

而$A_1A_2$的极点是$B_1$，$A_4A_5$极点是$B_4$，那么$A$的极线就是$B_1B_4$

同理$B$极线为$B_2B_5$，$C$为$B_3B_6$

由于$ABC$共线，那么对应的三线交点共点
\subsection{布利安桑证明帕斯卡}
如图，由配极，$A_iA_{i+1}$极点$B_i$

$B_1B_4$,$B_2B_5$,$B_3B_6$交于同一点

由上，$B_1B_4$极点就是$A$，$B_2B_5$为$B$，$B_3B_6$为$C$

那么ABC共线
\end{document}